%==============================================================================
\section{Discussion}\label{sec:disc}
%==============================================================================

\subsection{Bound-limited, not implementation-limited}\label{ssec:disc-bound}

A method can fail to solve an instance for two reasons. It can fail to find a good
solution, or it can find one and fail to prove that no better one exists. The
distinction matters because the two failures call for different remedies. The first
is answered by better search: heuristics, warm starts, branching rules. The second is
answered only by a better bound.

For the SSP the campaign answers this question unambiguously, and it answers it three
times over.

The first answer is in the direct measurement. On the runs that hit the time limit and
whose optimum is known from elsewhere, the Benders solver's incumbent already equals
that optimum $71\%$ of the time. It is not looking for the answer. It has the answer.

The second answer is in the ablation. Four strengthenings were implemented. The one
that improves the incumbent, a hybrid genetic search run at the root, gains
twenty-nine instances. The three that aim at the dual bound --- fractional cuts,
Pareto-optimal lifting, a conflict-graph row --- are neutral or harmful, and the worst
of them costs thirty-seven. A solver whose difficulty was finding good solutions would
not behave this way.

The third answer is the split by coverage-bound tightness. The same solver closes
$94\%$ of the instances on which the elementary bound $q$ happens to equal the
optimum, and $49\%$ of those on which it does not. That classification is not by size
or density. It is by a single structural property, and it is the property
Section~\ref{sec:gap} identifies as governing the whole problem.

The natural objection is that this is a statement about one implementation. It is not,
and Proposition~\ref{prop:pcf} is why. The position--configuration relaxation, arrived
at from a completely different direction, also equals $q$ exactly. So does the
multicommodity relaxation of \citet{daSilva2024}, proved by its own authors. Three
independent routes to the same value is not a coincidence of engineering; it is a
property of the problem. A configuration-based exact method for the SSP is limited by
the strength of its bound and not by its implementation.

The fractional-cut result deserves a separate remark, because it is the strongest form
of this conclusion available. Before the campaign, the honest position was that
fractional Benders cuts had never been tested here, because a defect had silently
discarded them. The temptation was to treat that as an open question with a promising
answer. The corrected run generated sixty million cuts and moved the bound on three
instances out of $1{,}238$. That is not an untested idea; it is a tested idea that
does not work on this problem, and knowing which of the two one has is the difference
between a research direction and a dead end.

\subsection{Where the theory and the measurements meet}\label{ssec:disc-agree}

Three agreements between Section~\ref{sec:gap} and Section~\ref{sec:comp} are worth
drawing out, because each was arrived at independently.

\paragraph{The bound governs both sides.}
Section~\ref{sec:gap} shows that $q$ equals the optimum on a large part of the
instance space, and that where it does the two-phase heuristic tends to be exact.
Section~\ref{sec:comp} shows that where it does, the exact methods close at the root,
and where it does not they branch without closing. The same quantity that makes the
heuristic exact makes the exact methods fast. That is not obvious in advance: it says
that the instances on which one would most want an exact method are exactly the
instances on which the available exact methods are weakest.

\paragraph{The gap opens where the magazine is large.}
Observation~\ref{obs:census} finds that every instance with a positive heuristic gap
has $b\ge4$, and that positive gaps are rare in a random sample. Independently,
\citet{Burger2015} report from an industrial study that the decomposition does well
when the largest job nearly fills the magazine, and that its solution quality
decreases as the gap between the magazine size and the largest job size grows. One of
these is an exhaustive combinatorial search on small instances and the other is a case
study on a printing press. They describe the same regime.

\paragraph{The walk model overstates the weakness of the standard method.}
This is the least expected of the three. The literature's mental model of the
two-phase heuristic is a walk through one magazine per group, and under that model the
canonical example family has a positive gap that grows without bound. Under the method
as it is actually specified, with its final loading step, that same family has no gap
at all (Observation~\ref{obs:rings}). The standard method is better than the standard
picture of it. Problem~\ref{prob:ktns-closes} asks how much better, and that question
did not exist before the two readings were separated.

\subsection{Limitations}\label{ssec:limits}

Several limitations bear on how far the conclusions above can be pushed.

\paragraph{The branch-and-price runs are not yet in.}
Their campaign under the uniform protocol is still running. The structural results
about their relaxations are unaffected, since those are proved rather than measured,
but the claim that a configuration branch-and-price is bound-limited currently rests on
Proposition~\ref{prop:pcf} and on the Benders measurements rather than on a direct
measurement of the prototypes.

\paragraph{Formulation 5 of \citet{Catanzaro2015} was not implemented.}
F4 was used as the compact baseline from that family. Their own conclusions recommend
F5 as the best of the family, and their experiments show it solving more instances
than F4. The comparison in Section~\ref{sec:comp} is therefore against the second
strongest member of that family, not the strongest, and the compact side of the
comparison is correspondingly favourable to the methods of this report.

\paragraph{The campaign is a snapshot.}
It had not finished when Section~\ref{sec:comp} was prepared, so not every
configuration has a result on every instance. Comparisons are made on common subsets
for this reason, and absolute solved counts across columns of Table~\ref{tab:solves}
are not comparable. The completed campaign will remove the caveat.

\paragraph{The prototypes are deliberately unaccelerated.}
The branch-and-price implementations carry no warm-started column pool, no dual
stabilisation and no heuristic pricing in the measured configuration. This is
intentional --- the prototype exists to measure the relaxation, and every
acceleration would confound that measurement --- but it means their solve counts are a
lower bound on what a tuned implementation would achieve.

\paragraph{The learned cut-selection prototype is still being evaluated.}
Section~\ref{ssec:rl} records what was built and on what it was validated. Its runs on
SSP instances are in progress, so no claim about its effect on this problem is made
here.

\paragraph{The setup-cost spectrum is not calibrated.}
Theorem~\ref{thm:collapse} locates the transition between the SSP and the JGP in terms
of the parameter $\rho$, but no source giving comparable measurements of machine stop
time and per-tool exchange time was found, so where real installations fall on that
spectrum is not established here. Remark~\ref{rem:rho} states this.
